\documentclass[11pt]{article}
\usepackage[margin=1.1in]{geometry}
\usepackage{fontspec}
\setmainfont{TeX Gyre Pagella}
\usepackage{amsmath}
\usepackage{enumitem}
\usepackage{hyperref}
\usepackage{xcolor}
\usepackage{fancyvrb}
\usepackage{titlesec}
\titleformat{\section}{\large\bfseries}{}{0em}{}
\titleformat{\subsection}{\normalsize\bfseries}{}{0em}{}
\setlength{\parskip}{0.4em}
\setlength{\parindent}{1.5em}

\title{\textbf{F\# Coding Standard for \textit{The Ledger of Meluhha}}\\
  \large Handoff --- Code Fix Plan}
\author{Third Buyer Advisory}
\date{Drafted April 2026 by the audit session on branch \texttt{claude/check-location-context-E8VFz}}

\begin{document}
\maketitle

%%% ===================================================================
\section{The one rule}

\textbf{The code is the paper.}

The audience for every \texttt{.fsx} script in this repository is the
reader of \textit{The Ledger of Meluhha} --- specifically, an
archaeology professor who has never written F\#.
The professor flips from paper \S N to script X, line Y, and must
recognise what the paper claims in what the code does.

This is the only rule.
Every other guideline below is a consequence.

Specifically:

\begin{itemize}[leftmargin=2em]
\item Names in code match vocabulary in the paper
  (\texttt{F1\_Merchant}, \texttt{F2\_Commodity}, not
  \texttt{field1}, \texttt{f2}).
\item Comments cite paper sections
  (\texttt{// Ledger of Meluhha §6.2 step 3}), not abstract
  engineering principles.
\item Data tables in code cite the paper table they mirror
  (\texttt{// Table 1, §2.2}).
\item The decode pipeline (\S6.2) reads as five named functions,
  ordered as the paper orders them.
\item Defensive code that manufactures unknowns or silently
  defaults is \emph{not} readable --- it hides the very
  decode-failures the research wants to surface.
\item Idiomatic F\# is a means, not an end.  Prefer legibility
  over type-level cleverness.
\end{itemize}

What the professor \emph{can} read: records, arrays, maps, pattern
matches on strings, SQL string literals, simple
\texttt{let}\,/\,\texttt{match} pipelines.
What the professor \emph{cannot} read without help: computation
expressions, type providers, active patterns, custom operators,
monads.
The first kind stays; the second kind is not introduced.

%%% ===================================================================
\section{Current state}

Five scripts implement the computational work:

\begin{itemize}[leftmargin=2em]
\item \texttt{indus\_seed\_codebook.fsx} --- seeds
  \texttt{indus\_codebook.db} from F\# literals.
  Mirrors \S6 Codebook.
\item \texttt{indus\_ingest\_corpus.fsx} --- ingests five sources
  into \texttt{indus\_corpus.db}.
  Mirrors \S6.1.
\item \texttt{indus\_decoder.fsx} --- demonstration decoder on
  hardcoded sign sequences.
  Not cited in \S7 Results.
\item \texttt{indus\_tn\_decoder.fsx} --- real CISI decode of TN
  parallels.
  Mirrors \S7.1.
\item \texttt{indus\_lssc.fsx} --- LSSC analysis.
  Mirrors \S6.3.
\end{itemize}

The code works and the numbers in the paper come from it.
What follows is readability-and-consistency work, not
bug-fixing.

%%% ===================================================================
\section{Priority-ordered fix list}

Fixes are grouped by how much risk running-without-a-compiler
carries.  P1 is safe to apply anywhere; P4 requires
\texttt{dotnet fsi} on a workstation.

\subsection{P1 --- blind-safe text edits (no compiler needed)}

\paragraph{P1.1 Rename \texttt{CargoTag} fields.}
In \texttt{indus\_decoder.fsx} lines 63--68, rename the record
fields so the record itself becomes Table~2 of the paper:

\begin{Verbatim}[fontsize=\small,frame=single]
type CargoTag = {
    F1_Merchant   : MerchantMark   // §4.1 Table 2
    F2_Commodity  : Commodity
    F3_WeightTier : WeightTier
    F4_Quantity   : int
    F5_Route      : Route }
\end{Verbatim}

Propagate the rename through \texttt{parseSeal} (the record
construction at \texttt{indus\_decoder.fsx:186--187}) and
\texttt{printTag} (\texttt{indus\_decoder.fsx:205--213}).

\paragraph{P1.2 Delete phantom header references.}
Every script header currently says:

\begin{Verbatim}[fontsize=\small]
// Conforms to: spec/fsharp/reference/fsharp_coding_standard.tex
\end{Verbatim}

Replace with a paper-section reference appropriate to the script:

\begin{Verbatim}[fontsize=\small]
// Implements: Ledger of Meluhha §6 (codebook seed)
// Coding standard: spec/fsharp/reference/fsharp_coding_standard.tex
\end{Verbatim}

(The standard file now exists at that path.  This is the one
pointing you read right now.)

\paragraph{P1.3 Cite paper tables in codebook literals.}
In \texttt{indus\_seed\_codebook.fsx} above each literal array,
add a comment pinning it to the paper:

\begin{Verbatim}[fontsize=\small]
// Table 1 (§2.2): Harappan weight progression.
// Base unit = 0.856 g [Hemmy 1931].  Binary series (×1…×64)
// for precious/indivisible goods; decimal (×160…×1000) for
// bulk cargo.
let weightTiers = [| ... |]
\end{Verbatim}

Do this for \texttt{weightTiers}, \texttt{commodities},
\texttt{routes}, \texttt{merchantMarks}, \texttt{signRoles},
\texttt{positionalRules}.

\paragraph{P1.4 Mark the demonstration decoder.}
\texttt{indus\_decoder.fsx} is not cited in \S7.1 (those results
come from \texttt{indus\_tn\_decoder.fsx}).  Either delete the
hardcoded seal data (lines 191--201), or prepend a clear banner:

\begin{Verbatim}[fontsize=\small]
// DEMONSTRATION only.  The results in §7.1 of the paper come
// from indus_tn_decoder.fsx over real CISI data.  This script
// exists to exercise the codebook on synthetic sign sequences.
\end{Verbatim}

\subsection{P2 --- structural alignment (compile-check required)}

\paragraph{P2.1 Extract the decode chain as a named pipeline.}
The paper (\S6.2) lists five steps:

\begin{enumerate}[leftmargin=2em]
\item Morphological parallel --- TN potsherd $\to$ IVC seal.
\item CISI inscription --- seal $\to$ Parpola sign sequence.
\item Concordance --- Parpola $\to$ Mahadevan.
\item Codebook --- Mahadevan $\to$ field role.
\item Lookup --- field role $\to$ content.
\end{enumerate}

Currently these five steps are inlined in the body of a
\texttt{for} loop at \texttt{indus\_tn\_decoder.fsx:248--293}.
Replace with a named function per step, composed in order:

\begin{Verbatim}[fontsize=\small,frame=single]
// Ledger of Meluhha §6.2: The decode chain
let decodeParallel (p: Parallel) : DecodedSeal =
    p
    |> findMorphologicalParallel   // §6.2 step 1
    |> loadCisiInscription         // §6.2 step 2
    |> mapParpolaToMahadevan       // §6.2 step 3
    |> lookupSignRole              // §6.2 step 4
    |> lookupFieldContent          // §6.2 step 5
\end{Verbatim}

The professor should be able to grep for ``step 3'' and find the
line.

\paragraph{P2.2 Eliminate stringly-typed roles.}
The SQL \texttt{CHECK(role IN
('commodity','weight','quantity','terminal','structural'))} is
a closed set.  In F\# it should be a DU parsed once at the SQL
read boundary:

\begin{Verbatim}[fontsize=\small,frame=single]
type Role =
    | Commodity  of CommodityCode
    | Weight     of WeightMultiplier
    | Quantity   of int
    | Terminal   of RouteCode
    | Structural
\end{Verbatim}

This replaces the string compares at
\texttt{indus\_decoder.fsx:160--178} and
\texttt{indus\_tn\_decoder.fsx:160--174}.

Same treatment for \texttt{Series}
(\texttt{binary}\,/\,\texttt{decimal}), \texttt{Position}
(\texttt{first}\,/\,\texttt{last}\,/\,\texttt{other}),
\texttt{FormType}, \texttt{Confidence}, \texttt{SignClass}.

\paragraph{P2.3 Raw SQL is banned on the read path.  Codebook is frozen typed records.}

\textbf{Rule.}  No \texttt{SqliteCommand}, \texttt{ExecuteReader},
or \texttt{GetString(n)} in any script that \emph{reads} codebook
data.  The codebook schema is frozen.  All read-path scripts use
\texttt{\#load "IndusCodebookTypes.fsx"} and typed F\# records
with prebuilt lookup maps.  No SQL at runtime.  No database
connection for the codebook.

\textbf{How it works.}
SqlHydra\footnote{\url{https://github.com/JordanMarr/SqlHydra}}
v4 was run \emph{once} against \texttt{indus\_codebook.db} to
generate typed records matching the schema
(\texttt{hydra/IndusCodebook.fs}).  Those records were then frozen
--- with literal data arrays --- into
\texttt{IndusCodebookTypes.fsx}.  The generator was discarded.
No \texttt{SqlHydra.Query} runtime dependency exists.

\begin{Verbatim}[fontsize=\small,frame=single]
#load "IndusCodebookTypes.fsx"
open IndusCodebookTypes

// Typed record, prebuilt map, zero DB:
let jar = commodityByCode |> Map.find "jar"
// jar.Label = "JAR GOODS", jar.Description = "sesame oil/grain/resin"
\end{Verbatim}

\texttt{jar.Label} is a typed string field, not
\texttt{r.GetString 0}.  Rename a field $\to$ compile error.
The professor opens \texttt{IndusCodebookTypes.fsx}, reads
28~sign roles, 8~commodities, 5~routes as plain F\# records.
The schema \emph{is} the documentation.

\textbf{Two exceptions} (write-path scripts that create databases):

\begin{enumerate}[nosep]
\item \texttt{indus\_seed\_codebook.fsx} --- DDL +
  \texttt{INSERT}.  Creates the codebook.  Raw ADO.NET is the
  only option for schema creation.
\item \texttt{indus\_ingest\_corpus.fsx} --- ETL pipeline, bulk
  inserts from external sources into \texttt{indus\_corpus.db}.
  Same reason.
\end{enumerate}

\noindent
Both are write-path.  Every script that reads codebook data is
banned from raw SQL.  This is enforced by the pre-commit hook
(coding standard audit on staged \texttt{.fsx} files).

\textbf{When the codebook changes} (rare --- stable at 28~signs,
8~commodities, 5~routes since initial seeding):

\begin{Verbatim}[fontsize=\small,frame=single]
# 1. Edit indus_seed_codebook.fsx (source of truth)
# 2. cmake --build build --target codebook
# 3. cd hydra && sqlhydra sqlite   (re-generate .fs)
# 4. Update IndusCodebookTypes.fsx (frozen data + types)
\end{Verbatim}

\subsection{P3 --- behavioural changes (run decoder before/after)}

\paragraph{P3.1 Remove defensive unknowns.}
\texttt{indus\_decoder.fsx:133--136} manufactures
\texttt{unknownCommodity}, \texttt{unknownRoute},
\texttt{unknownMerchant} so \texttt{parseSeal} always returns a
\texttt{CargoTag}.  Replace with a DU that distinguishes success
from partial resolution:

\begin{Verbatim}[fontsize=\small,frame=single]
type Decoded =
    | CargoTag of F1_Merchant * F2_Commodity * F3_WeightTier
                * F4_Quantity * F5_Route
    | OwnershipMarkOnly of MerchantMark
    | PartiallyResolved of MerchantMark
                         * (Role * Resolution) list
\end{Verbatim}

A partial decode is a research output (§7.4: ``Expanding the
codebook is the next phase''), not glue to hide.

Run before: \texttt{dotnet fsi indus\_tn\_decoder.fsx
indus\_corpus.db indus\_codebook.db > before.txt}.
Run after, \texttt{diff before.txt after.txt} --- the decoded
M-52A and M-148A cargo contents must be identical.  Only the
unresolved-sign reporting format should differ.

\paragraph{P3.2 Verify M-52A and M-148A decode.}
The paper's current \S7.1 says:

\begin{itemize}[leftmargin=2em]
\item M-52A $\to$ jar goods / Mesopotamia
\item M-148A $\to$ jar goods / domestic
\end{itemize}

After P1 and P2 are applied, rerun the decoder and confirm this
output.  If it differs, either the paper or the code is wrong
--- flag and stop.

\subsection{P4 --- schema fixes (touches ingest + reader)}

\paragraph{P4.1 JSON blob $\to$ join table.}
\texttt{cisi\_inscription.signs TEXT} and
\texttt{sign\_concordance.mahadevan\_ids TEXT} store collections
as JSON-ish strings.  Alloy specifies them as \texttt{seq
BaseSign} and \texttt{set BaseSign}.  The corresponding F\#
readers hand-parse with \texttt{.Trim('[', ']').Split(',')}
(\texttt{indus\_tn\_decoder.fsx:100--104, 116--119}) --- a
parser for a format with no grammar.

Add two join tables in \texttt{indus\_ingest\_corpus.fsx}:

\begin{Verbatim}[fontsize=\small,frame=single]
CREATE TABLE cisi_inscription_sign (
    inscription_id TEXT NOT NULL
        REFERENCES cisi_inscription(id),
    position       INTEGER NOT NULL,
    parpola_id     TEXT NOT NULL,
    PRIMARY KEY (inscription_id, position)
) STRICT;

CREATE TABLE sign_concordance_mahadevan (
    parpola_id      TEXT NOT NULL
        REFERENCES sign_concordance(parpola_id),
    mahadevan_id    TEXT NOT NULL,
    PRIMARY KEY (parpola_id, mahadevan_id)
) STRICT;
\end{Verbatim}

Keep the JSON TEXT columns for one release (migration path) or
drop them immediately and update the readers.  Alloy checks A10
(concordance crosses sources) and A12 (CISI signs valid) become
enforceable.

\paragraph{P4.2 Add the \texttt{artefact} table.}
Alloy \S A1/A2 assert \texttt{Artefact.site: one Site} and
\texttt{SignOccurrence.artefact: one Artefact}.  Currently
\texttt{sign\_occurrence.artefact\_id} is an orphan TEXT ---
no target table exists.  Add:

\begin{Verbatim}[fontsize=\small,frame=single]
CREATE TABLE artefact (
    id          TEXT PRIMARY KEY,
    site_id     TEXT NOT NULL REFERENCES site(id),
    source_code TEXT NOT NULL REFERENCES source(code)
) STRICT;
\end{Verbatim}

Make \texttt{sign\_occurrence.artefact\_id} an FK to it.

\paragraph{P4.3 Preserve iron-DB invariants on re-projection.}
\texttt{indus\_ingest\_corpus.fsx:168--202} re-emits
\texttt{radiometric\_site}, \texttt{radiometric\_sample},
\texttt{radiometric\_measurement} into the corpus DB with all
\texttt{CHECK} constraints stripped (P1/P2/P4/P5/P9 from
\texttt{data/audit\_iron\_db.als}).  Two options:

\begin{enumerate}[leftmargin=2em]
\item Re-declare the same CHECKs on the target tables.
\item Replace the copy with \texttt{ATTACH DATABASE
  'antiquity\_of\_iron.db' AS iron} and query through the attach,
  so the frozen DSL is never duplicated.
\end{enumerate}

Option 2 is cleaner but changes the deployment surface.  Prefer
1 unless attach is acceptable.

\paragraph{P4.4 Add \texttt{UNIQUE(sample, institution, method)}}
on \texttt{radiometric\_measurement}, matching A12 of the iron
DSL.  Promote \texttt{institution} from TEXT to an FK to an
\texttt{institution} table.

\paragraph{P4.5 Write \texttt{spec/indus\_codebook.als}.}
The codebook has real invariants (a sign has one role; a
\texttt{commodity}-role sign has a \texttt{ref\_code} but no
\texttt{ref\_multiplier}; a \texttt{weight}-role sign has a
\texttt{ref\_multiplier} but no \texttt{ref\_code}).  Today
these are enforced only by the CHECK IN string on
\texttt{role}.  Add an Alloy spec and bring codebook
verification up to the iron-DB tier.

%%% ===================================================================
\section{Verification protocol}

Before and after \emph{every} P2, P3, or P4 change:

\begin{enumerate}[leftmargin=2em]
\item \texttt{dotnet fsi indus\_seed\_codebook.fsx indus\_codebook.db}
\item \texttt{dotnet fsi indus\_ingest\_corpus.fsx
  ../indus-valley indus\_corpus.db}
\item \texttt{dotnet fsi indus\_tn\_decoder.fsx
  indus\_corpus.db indus\_codebook.db > decode.txt}
\item Confirm \texttt{decode.txt} contains M-52A $\to$ jar
  goods / Mesopotamia and M-148A $\to$ jar goods / domestic.
\item \texttt{dotnet fsi indus\_lssc.fsx <corpus.csv> lssc.db}
  and confirm the paper's \S6.3 numbers: 7{,}637 tokens,
  3{,}726 unique signs, 7{,}037 transitions, LSSC mean 0.69.
\end{enumerate}

If any of these shift after a P2/P3/P4 change, stop and
reconcile before continuing.

For Alloy (P4.5 and the count reconciliation flagged in the
paper):

\begin{Verbatim}[fontsize=\small]
java -jar alloy.jar spec/indus_corpus.als
java -jar alloy.jar data/audit_iron_db.als
\end{Verbatim}

The paper's abstract, \S6.1, and \S A.3 claim twelve assertions
UNSAT at scope~6.  The current spec has sixteen.  Either
re-run and update the count (to sixteen) or prune the spec back
to twelve --- probably by moving the four radiometric
assertions (A13--A16) into a separate
\texttt{spec/indus\_radiometric.als} since they are really
about the iron-DB subsumption, not the corpus qua corpus.

%%% ===================================================================
\section{What's already done (on branch \texttt{claude/check-location-context-E8VFz})}

Commit \texttt{5089ffe} reframes the paper:

\begin{itemize}[leftmargin=2em]
\item Abstract, \S1, \S7, \S9 rewritten to name the 100-year
  impasse and position the paper as Path~C between Path~A
  (linguistic decipherment, $\sim$100 failed attempts) and Path~B
  (Farmer--Sproat--Witzel 2004).
\item \S8 and \S11 cleaned up: the phantom M-67A reference
  removed; M-52A and M-148A decodes made consistent with \S7.1
  (both jar goods; M-52A $\to$ Mesopotamia; M-148A $\to$
  domestic).
\item Parpola 1994 and Parpola 2008 added to
  \texttt{citations.lua}.
\item TODO markers placed at the three locations that claim
  twelve Alloy checks, pending a re-run on a machine with
  \texttt{java} and \texttt{alloy.jar}.
\end{itemize}

No code changes on this branch yet.  The code fixes P1--P4
above are the next session's work.

%%% ===================================================================
\section{Handoff notes}

This document lives at
\texttt{spec/fsharp/reference/fsharp\_coding\_standard.tex},
which is exactly the path every \texttt{.fsx} script header
cites.  Until now that path was vapour.  It now exists.

The workstation Claude with \texttt{dotnet fsi} available
should:

\begin{enumerate}[leftmargin=2em]
\item Read this file.
\item Read the paper (\texttt{ledger\_of\_meluhha.tex}) ---
  especially \S1, \S4, \S6, \S7.
\item Apply P1 in one commit (mechanical, low risk).
\item Apply P2 in one commit per sub-item, with compile-check
  after each.
\item Apply P3 with the before/after \texttt{diff} protocol.
\item Apply P4 last, one sub-item per commit, each gated on the
  verification protocol of \S4 above.
\item When P4 is done and all 16 Alloy checks are UNSAT, update
  the paper's count claims (abstract, \S6.1, \S A.3) and remove
  the TODO comments.
\end{enumerate}

Do not introduce type providers, active patterns, custom
operators, or generic computation expressions (\texttt{maybe
\{ \}}, custom builders).  The professor reads the code.  She
does not read Wlaschin.

One exception: \texttt{query \{ \}} is F\#'s native SQL DSL.
It reads exactly like SQL because it \emph{is} SQL, spelled in
F\#.  A reader who can read SQL can read a query expression.
It stays.  See P2.3.

\end{document}
